\chapter{محاسبه با عملگرها}%
\label{ch:operators}

برنامه‌ها دائم در حالِ حساب کردن‌اند: جمعِ خرید، میانگینِ نمره، مقایسهٔ
امتیازها. علامت‌هایی مانندِ \code{+} و \code{*} که این کارها را انجام
می‌دهند \textbf{عملگر} نام دارند. در این فصل سه خانوادهٔ عملگرها را
می‌شناسیم: حسابی، مقایسه‌ای و منطقی.

\section{عملگرهای حسابی}

\begin{center}
\begin{tabular}{clc}
\toprule
عملگر & کار & نمونه و حاصل \\
\midrule
\code{+} & جمع              & \code{7 + 2} $\rightarrow$ \code{9} \\
\code{-} & تفریق            & \code{7 - 2} $\rightarrow$ \code{5} \\
\code{*} & ضرب              & \code{7 * 2} $\rightarrow$ \code{14} \\
\code{/} & تقسیم            & \code{7 / 2} $\rightarrow$ \code{3} \\
\code{\%} & باقی‌مانده      & \code{7 \% 2} $\rightarrow$ \code{1} \\
\bottomrule
\end{tabular}
\end{center}

جمع و تفریق و ضرب همان‌اند که در دبستان آموختیم، اما دو سطرِ آخرِ جدول
داستان دارند.

\subsection{تقسیمِ صحیح: خارج‌قسمت، بدونِ اعشار}

حاصلِ \code{7 / 2} در جدول \code{3} نوشته شده، نه \code{3.5}. اشتباه
چاپی نیست! وقتی هر دو طرفِ تقسیم عددِ \emph{صحیح} باشند، حاصل هم صحیح است:
فقط خارج‌قسمت را نگه می‌دارد و اعشار را دور می‌ریزد. اگر جوابِ اعشاری
می‌خواهید، دستِ‌کم یکی از دو طرف باید اعشاری باشد:

\begin{salam}
تابع اصلی:
    چاپ 7 / 2
    چاپ 7.0 / 2.0
تمام
\end{salam}

\begin{progout}
3
3.5
\end{progout}

\begin{warn}
این رفتار یکی از پرتکرارترین دام‌های تازه‌کارهاست. اگر میانگینِ ۷ و ۸ را با
\code{(7 + 8) / 2} حساب کنید، به‌جای \code{7.5} عددِ \code{7} را می‌گیرید و
شاید ساعت‌ها دنبالِ اشکال بگردید. هر جا میانگین یا نسبت حساب می‌کنید، به
صحیح یا اعشاری بودنِ عملوندها فکر کنید.

دامِ دوم: تقسیم بر صفر مجاز نیست و برنامه را با خطا متوقف می‌کند. یادتان
هست تمرینِ فصلِ \ref{ch:algorithms} پرسید «اگر کلاس هیچ نمره‌ای نداشته
باشد چه می‌شود؟» پاسخ همین است: شمارنده صفر می‌ماند و تقسیم بر صفر برنامه
را از پا درمی‌آورد. الگوریتمِ خوب پیش از تقسیم، صفر نبودنِ مخرج را بررسی
می‌کند؛ ابزارش را در فصلِ بعد می‌سازیم.
\end{warn}

\subsection{باقی‌مانده: عملگرِ کم‌ادعا و پرکاربرد}

\code{\%} باقی‌ماندهٔ تقسیم را می‌دهد: \code{17 \% 5} می‌شود \code{2}، چون
۱۷ سه تا ۵ دارد و ۲ تا باقی می‌ماند. به ظاهرِ ساده‌اش نگاه نکنید؛ این عملگر
همه‌جا هست:

\begin{itemize}
  \item \textbf{زوج یا فرد؟} عددی زوج است اگر و فقط اگر
        \code{عدد \% 2} برابرِ \code{0} باشد.
  \item \textbf{بخش‌پذیری:} \code{عدد \% 7 == 0} یعنی عدد مضربِ ۷ است.
  \item \textbf{رقمِ یکان:} \code{عدد \% 10} رقمِ یکانِ هر عدد را می‌دهد؛
        مثلاً \code{1404 \% 10} می‌شود \code{4}.
\end{itemize}

\section{اولویت و پرانتز}

عبارتِ \code{2 + 3 * 4} چند می‌شود؟ اگر از چپ حساب کنید ۲۰، اما جوابِ درست
\code{14} است: ضرب و تقسیم \emph{اولویتِ} بالاتری از جمع و تفریق دارند،
درست مانندِ ریاضیِ مدرسه. هر جا ترتیبِ دیگری می‌خواهید (یا حتی فقط برای
خوانایی) از پرانتز استفاده کنید:

\begin{salam}
تابع اصلی:
    چاپ 2 + 3 * 4
    چاپ (2 + 3) * 4
تمام
\end{salam}

\begin{progout}
14
20
\end{progout}

\begin{tip}
وقتی شک دارید، پرانتز بگذارید. پرانتزِ اضافه هیچ هزینه‌ای ندارد اما ابهام را
هم برای شما و هم برای خواننده از بین می‌برد.
\end{tip}

\section{عملگرهای مقایسه‌ای}

خانوادهٔ دوم، دو مقدار را می‌سنجد و حاصلش همیشه یک مقدارِ \code{منطقی}
است: \code{درست} یا \code{نادرست}.

\begin{center}
\begin{tabular}{clc}
\toprule
عملگر & معنی & نمونه و حاصل \\
\midrule
\code{==} & برابر است؟          & \code{5 == 5} $\rightarrow$ \code{درست} \\
\code{!=} & نابرابر است؟        & \code{5 != 5} $\rightarrow$ \code{نادرست} \\
\code{<}  & کوچک‌تر است؟        & \code{3 < 7} $\rightarrow$ \code{درست} \\
\code{>}  & بزرگ‌تر است؟        & \code{3 > 7} $\rightarrow$ \code{نادرست} \\
\code{<=} & کوچک‌تر یا مساوی؟   & \code{7 <= 7} $\rightarrow$ \code{درست} \\
\code{>=} & بزرگ‌تر یا مساوی؟   & \code{3 >= 7} $\rightarrow$ \code{نادرست} \\
\bottomrule
\end{tabular}
\end{center}

\begin{salam}
تابع اصلی:
    سن : صحیح = 15
    چاپ سن >= 18
    چاپ سن < 18
تمام
\end{salam}

\begin{progout}
false
true
\end{progout}

\begin{warn}
مهم‌ترین هشدارِ این فصل: \code{=} و \code{==} دو چیزِ کاملاً متفاوت‌اند.
\code{=} یعنی «بگذار در جعبه» (انتساب) و \code{==} یعنی «آیا برابرند؟»
(مقایسه). تقریباً همهٔ تازه‌کارها دستِ‌کم یک بار این دو را جابه‌جا
می‌نویسند؛ شما هم نوشتید، نگران نشوید؛ فقط یادتان باشد اولین جایی که
باید دنبالِ اشکال بگردید همین‌جاست.
\end{warn}

\section{عملگرهای منطقی}

خانوادهٔ سوم، مقدارهای منطقی را با هم ترکیب می‌کند. سه عضو دارد:

\begin{itemize}
  \item \code{\&\&} (\textbf{و}): وقتی درست است که \emph{هر دو} طرف درست
        باشند.
  \item \code{||} (\textbf{یا}): وقتی درست است که \emph{دستِ‌کم یکی} از دو
        طرف درست باشد.
  \item \code{!} (\textbf{نقیض}): درست را نادرست می‌کند و برعکس.
\end{itemize}

با این‌ها می‌توان شرط‌های مرکبِ دنیای واقعی را نوشت. «برای این بازی باید
دستِ‌کم ۷ سال داشته باشی و قدت از ۱۲۰ بیشتر باشد»:

\begin{salam}
تابع اصلی:
    سن : صحیح = 9
    قد : صحیح = 130
    مجاز : منطقی = سن >= 7 && قد > 120
    چاپ "اجازه دارد؟", مجاز
    چاپ "اجازه ندارد؟", !مجاز
تمام
\end{salam}

\begin{progoutfa}
اجازه دارد؟ true
اجازه ندارد؟ false
\end{progoutfa}

\begin{note}
«یا»ی برنامه‌نویسی با «یا»ی محاوره کمی فرق دارد: در گفتار، «چای یا قهوه؟»
یعنی فقط یکی؛ اما \code{||} وقتی هر دو طرف هم درست باشند، درست است. به آن
«یای فراگیر» می‌گویند.
\end{note}

\section{الگوریتم‌ها بالاخره به کد می‌رسند}

حالا که حساب و مقایسه داریم، الگوریتمِ «زوج یا فرد» از تمرین‌های فصلِ
\ref{ch:algorithms} را تا نیمه پیاده کنیم؛ نتیجهٔ بررسی را فعلاً به صورتِ
\code{درست}/\code{نادرست} چاپ می‌کنیم و در فصلِ بعد که «انتخاب» را یاد
گرفتیم، پیامِ فارسیِ زیبا هم می‌دهیم:

\begin{salam}
تابع اصلی:
    عدد : صحیح = 1404
    زوج است : منطقی = عدد % 2 == 0
    چاپ "عدد:", عدد
    چاپ "زوج است؟", زوج است
تمام
\end{salam}

\begin{progoutfa}
عدد: 1404
زوج است؟ true
\end{progoutfa}

\section{تمرین‌ها}

\exercise حاصلِ هر عبارت را اول روی کاغذ حساب کنید، سپس با \code{چاپ}
بررسی کنید: \code{10 - 4 * 2}؛ \code{(10 - 4) * 2}؛ \code{9 / 2}؛
\code{9.0 / 2.0}؛ \code{9 \% 2}.

\exercise دمای هوا ۲۳ درجه است. عبارتی منطقی بنویسید که «هوای مطبوع» را
نشان دهد: دما دستِ‌کم ۱۸ و حداکثر ۲۸ باشد. حاصل را چاپ کنید.

\exercise برنامه‌ای بنویسید که برای عددِ \code{2358}، رقمِ یکان و رقمِ
دهگانش را جدا چاپ کند (راهنمایی: یکان با \code{\% 10} به دست می‌آید؛ برای
دهگان اول با \code{/ 10} یک رقم را بیندازید).

\exercise سال‌های کبیسهٔ میلادی: سالی کبیسه است که بر ۴ بخش‌پذیر باشد،
مگر آنکه بر ۱۰۰ بخش‌پذیر باشد؛ ولی اگر بر ۴۰۰ بخش‌پذیر باشد باز کبیسه
است. عبارتی منطقی برای «کبیسه بودنِ» متغیرِ \code{سال} بنویسید و آن را با
سال‌های 2024 (کبیسه)، 1900 (نیست) و 2000 (کبیسه) بیازمایید.

\exercise بدونِ اجرا بگویید حاصلِ \code{!(3 > 5) \&\& (2 == 2)} چیست؛
سپس با برنامه بررسی کنید.
